\section{Faza I — wyłonienie reprezentanta każdego wariantu}
\label{sec:phase1}

Celem fazy I było wyznaczenie jednego reprezentanta dla każdego z sześciu badanych wariantów agenta. Reprezentant to wektor wag, który uzyskał dobrą jakość w przebiegu optymalizacji i potwierdził ją w procedurze porównawczej opisanej w sekcji~\ref{sec:analysis-scheme}.

Ta selekcja jest istotna z metodologicznego punktu widzenia. W zadaniach koewolucyjnych wynik pojedynczego przebiegu zależy częściowo od składu populacji przeciwników i losowości symulacji. Dwuetapowy turniej --- najpierw wewnątrz każdego optymalizatora, a potem między optymalizatorami --- zmniejsza ryzyko wyboru osobnika, który był dobry przypadkowo.

Wyjątek dotyczy wariantu \textit{21 głębokościowego}. Ze względu na około \(3{,}5\)-krotnie dłuższy czas pojedynczego treningu, uruchamiano po jednym agencie na każdego optymalizatora, dlatego dla tego wariantu pominięto pierwszy mini-turniej wewnątrz optymalizatora.

\begin{table}[ht]
\centering
\small
\caption{Reprezentanci wyłonieni w fazie I oraz odpowiadający im zwycięski optymalizator.}
\label{tab:phase1-representatives}
\begin{tabular}{ll}
\toprule
Wariant & Schemat optymalizacji wybranego reprezentanta \\
\midrule
\textit{bazowy 21}          & SHADE koewolucyjny \\
\textit{28 parametrów}      & SHADE koewolucyjny \\
\textit{28 znormalizowany}  & SHADE koewolucyjny \\
\textit{63 parametry}       & SHADE koewolucyjny \\
\textit{63 płynny}          & SHADE koewolucyjny \\
\textit{21 głębokościowy}   & SHADE z Hall of Fame \\
\bottomrule
\end{tabular}
\end{table}

Tabela \ref{tab:phase1-representatives} pokazuje dwie rzeczy. SHADE koewolucyjny był najskuteczniejszy dla pięciu z sześciu wariantów, więc wygląda na dobrze dopasowanego do większości badanych konfiguracji. Z kolei dla wariantu głębokościowego lepszy wynik uzyskał SHADE z Hall of Fame --- najpewniej dlatego, że bardziej zróżnicowany zbiór przeciwników ma tu większe znaczenie.

Na tym etapie nie rozstrzyga się jeszcze, który wariant agenta jest najlepszy --- faza I służyła wyłącznie do wyboru reprezentantów. Widać jednak, że o jakości agenta decyduje również sposób prowadzenia optymalizacji.

\paragraph{Wniosek fazy I --- odpowiedź na RQ4.}
Algorytm koewolucyjny nie wyłonił najlepszego reprezentanta w żadnym z badanych wariantów, ponieważ we wszystkich sześciu przypadkach zwyciężył jeden z wariantów SHADE zaproponowanych w tej pracy. SHADE koewolucyjny był dominujący (5 z 6 wariantów). Odpowiedź na pytanie badawcze nr 4 (RQ4) jest więc jednoznaczna: oba warianty SHADE przewyższają algorytm koewolucyjny przy równoważnym budżecie obliczeniowym, choć wybór między nimi zależy od rodzaju modyfikacji agenta.
